% Pursuing Stacks (cotes 134-2 to 134-8): archive cote 134-6, site page 63 (PDF page 64), typescript page 381.
% typescript. Draft compared: ps5.tex:4557-4628.
% AI-transcribed by Claude Opus 5.5 (claude-opus-5-5, adaptive thinking, effort high), Message Batches API, batch ps3 (msgbatch_01XdFmuVVBAcgc7EFxdtCqsj), 2026-11-12; reviewed against both supplied scan halves; see CHECK_DECISIONS.tsv.
\scanpage{134-6/63}
\mspage{381}

(11) \[ B^{\wedge}_{ab} \hand{\xrightarrow{\approx}} \underline{\text{Hom}}_{!}(A^{\wedge}_{ab}, (\text{Ab})) \quad , \]
where the index ! denotes the full subcategory of \uline{Hom} made up with
functors commuting to small direct limits. It is immediate that projecti
% NOTE: the right edge of the scan cuts the end of the line above ("projecti...")
-ve objects in $B^{\wedge}_{ab}$ give rise to \xout{exact} objects in $\underline{\text{Hom}}_{!}$ which are \hand{\uline{\textcolor{black}{exact}}}
functors from $A^{\wedge}_{ab}$ to\deleted{\xout{f}} (Ab), and I'll have to check that the converse
also holds. I\over{f}{f} so, a \xout{projective} chain complex in $B^{\wedge}_{ab}$ with projective
components can be interpreted as being just an \xout{exact} \hand{\uline{\textcolor{black}{arbitrary}}} \hand{\uline{\textcolor{black}{exact}}}
\hand{\uline{\textcolor{black}{functor}}} \hand{\uline{\textcolor{black}{commuting to small sums}}}
\[ A^{\wedge}_{ab} \hand{\longrightarrow} \text{\over{C}{C}h.}(\text{Ab}) \quad , \]
\hand{(}never minding whether or not it can be described ``computationally'' in
terms of objects in Add(B) or in Addinf\over{f}{(}B)\unsure{\hand{!}}\hand{)} - which is all that is nee-%
ded in order to complete the diagram (8), and wonder if it commutes up
to isomorphism ! And the most natural statement here is that this is
indeed so whenever this functor, viewed as a chain complex in the \inserted[typed]{abelian} cate-%
gory
(12) \[ \underline{\text{Hom}}_{!}(A^{\wedge}_{ab},(\text{Ab})) \quad , \]
is a \inserted[typed]{(projective)} resolution of the canonical object $\overset{\hand{\sim}}{\underline{Z}}_{B}$ of the category (12), coming
% NOTE: the right edge of the scan cuts the end of the line above ("coming")
from the object $\underline{Z}_{B}$ of the left hand side of (11). Now this functor is
just the \xout{functor} familiar ``direct limit'' functor
(13) \[ \overset{\hand{\sim}}{\underline{Z}}_{B} \hand{\simeq} \underset{\over[typed]{B}{B}}{\underset{\hand{\longrightarrow}}{\text{lim}}} \quad : \quad A^{\over{a}{\wedge}}_{ab} \overset{\hand{\text{def}}}{=} \xout{\unsure{\text{Hom}(\text{Ab})}},\underline{\text{Hom}}(B,(\text{Ab})) \hand{\longrightarrow} (\text{Ab}) \quad , \]
which can be equally interpreted as \xout{\illegible{1}}
% This unresolved fragment occupies two overlapping lines of typewriter x strikes.
(14) \[ \overset{\hand{\sim}}{\underline{Z}}_{B} \hand{\simeq} p^{ab}_{A!} \quad : \quad A^{\wedge}_{ab} \hand{\longrightarrow} (\text{Ab}) \quad , \]
namely \inserted[typed]{(abelian)} ``integration'' with re\over[typed]{p}{s}\over[typed]{s}{p}ect to the map in (Cat)
\[ p_{A} \ : \ A \hand{\longrightarrow} \hand{\Delta}_{o} \ . \]
Thus, ulti\over[typed]{p}{m}ately, \hand{\uline{\textcolor{black}{abelianizators for}}} A (or what we may call ``sta\over[typed]{m}{n}dard
abelia\over[typed]{m}{n}izators'' , if there should turn out to be any others, and that
they are worth looking at ) \hand{\uline{\textcolor{black}{turn out to be no more,}}} \hand{\uline{\textcolor{black}{no less than just}}}
\hand{\uline{\textcolor{black}{a projective resolution,}}} \hand{\uline{\textcolor{black}{in the category}}} (12) \hand{\uline{\textcolor{black}{of functors from}}} $A^{\wedge}_{ab}$ \hand{\uline{\textcolor{black}{to}}}
(Ab) \hand{\uline{\textcolor{black}{commuting with small direct limits,}}} \hand{\uline{\textcolor{black}{of the most interesting object}}}
% NOTE: the right edge of the scan cuts the end of the line above ("object")
\hand{\uline{\textcolor{black}{in this category,}}} \hand{\uline{\textcolor{black}{namely the functor}}}
(15) \[ p^{ab}_{A!} \hand{\simeq} \underset{B}{\underset{\hand{\longrightarrow}}{\text{lim}}} \quad : \quad A^{\wedge}_{ab} \hand{\longrightarrow} (\text{Ab}) \quad . \]
We are far indeed from the faltering reflexions of section 91\hand{,} about
computing homology and cohomology \xout{in terms} of \xout{models} homotopy models
described \xout{by} in terms of test categories deduced some way or other from
cellular decompositions of spheres !

\margin{63}
